\subsubsection{The invoice and the work need different clocks}
\label{sec:mixed-contract-billing}

The builder can earn the same construction and inspection revenue while being
paid on a different schedule. Retain the preceding allocations, costs and
approved amendment, but place all work and payments in calendar 2024. At
inception and after the amendment, the expected interval between each payment
and its related service is at most one year, and the builder elects the
financing expedient. Receivables and contract assets have a separately assessed
zero impairment allowance in this example. These findings matter: AASB~15
requires the financing assessment and AASB~9 treatment of the rights to
consideration; eventual collection alone establishes neither
\citep[paras.~63,105--109]{aasb15dated}.

Suppose the customer initially advances only 30. Opening cash and capital of
100 become cash 130 and a contract liability of 30. In April the builder pays
construction costs of 40 and visit costs of 2, earning the same revenue of
55 as before. Cash falls to 88. The advance covers the first 30 of that work;
the remaining 25 is a contract asset because the right to payment still
depends on a billing milestone. Profit is 13 and equity is 113.

Now the milestone makes a gross invoice of 80 unconditional. The invoice
credits the advance of 30 already received, leaving a receivable of 50. The
contract asset disappears. But the builder has earned only 55 of the 80 now
paid or due, so a contract liability of 25 also appears. Figure
\ref{fig:mixed-billing-clocks} isolates this change. The receivable and contract
liability coexist: one is an unconditional payment right and the other is
work still owed under the same contract.

\begin{figure}[htbp]
\centering
\begin{tikzpicture}[>=Latex, every node/.style={font=\small},
  stage/.style={draw,rounded corners,align=center,text width=3.25cm,
  minimum height=1.65cm,inner sep=5pt}]
\node[stage] (advance) {Advance 30\\Cash 130\\Contract liability 30};
\node[stage,right=6mm of advance] (work) {Earn 55; pay costs 42\\Cash 88\\Contract asset 25};
\node[stage,right=6mm of work] (invoice) {Invoice 80; credit 30\\Receivable 50\\Contract liability 25};
\draw[->] (advance)--(work);
\draw[->] (work)--(invoice);
\end{tikzpicture}
\caption{The invoice changes the payment right without new work or new cash.
Cash stays at 88 and cumulative revenue at 55 when the final box is reached.}
\label{fig:mixed-billing-clocks}
\end{figure}

The advance must therefore retain its identity after receipt. Crediting it
reduces the new amount collectible; it does not refund cash or recognise
revenue again. Let \(G\) be cumulative gross invoices, \(K\) the advances
credited against those invoices, and \(C\) later collections of the resulting
receivables. The amount still collectible is
\[
  \mathrm{AR}=G-K-C.
\]
For the invoice just issued, this is \(80-30-0=50\). Collecting 20 changes
cash to 108 and the receivable to 30, with no change in revenue. Each credit
names both its advance and invoice; each collection names its invoice. A
second credit cannot reuse an advance that has already been consumed.

To determine the contract balance, count consideration already paid or
unconditionally due once. If \(A\) is cumulative advance cash, adding it to
gross invoices would count credited advances twice. Subtract \(K\) to obtain
\(T=G+A-K\). Compare this with cumulative earned revenue \(R\):
\[
  T=G+A-K,\qquad
  \mathrm{CA}=\max(R-T,0),\qquad
  \mathrm{CL}=\max(T-R,0).
\]
At invoicing, \(T=80+30-30=80\), giving the liability \(80-55=25\).
Collection leaves \(T\) unchanged: cash replaces a right that was already
unconditional. Together these relations imply, for cumulative paid costs
\(E\) and opening capital 100,
\[
  \text{cash}+\mathrm{AR}+\mathrm{CA}-\mathrm{CL}
       =100+R-E.
\]
Substitution shows why: cash is \(100+A+C-E\), the receivable is
\(G-K-C\), and the net contract asset is \(R-G-A+K\). The advance,
invoice, credit and collection terms cancel, leaving earned profit in equity.
The cancellation checks the arithmetic; the contractual findings determine
which amounts belong in it.

The May amendment now adds construction revenue of 6 without a receipt,
reducing the liability from 25 to 19. A new, separately identified advance of
50 raises cash to 158 and the liability to 69. Two further visits earn 10
and cost 4, leaving cash 154, receivables 30, contract liability 59 and profit
25. The remaining gross invoice of 90 credits the new advance of 50. Its net
right of 40 increases total receivables to 70; paid-or-due consideration
reaches 170, so the contract liability becomes \(170-71=99\).

Collect the first invoice's remaining 30, then complete construction and the
last three visits for cash costs of 66. Completion earns the remaining 99 and
removes the contract liability. Cash is 118, the second invoice still carries
40, and equity is 158. Collecting that 40 leaves cash and equity of 158, as
in the prepaid example. Customer cash is advances \(30+50\) plus collections
\(20+30+40\), totalling 170. Paid costs are 112 and profit is 58. The gross
invoices of 170 describe the rights created, not an additional receipt.

Full prepayment provides a useful limiting case. If the final invoice credits
all 170 already received, it creates no receivable and no monetary journal;
it still records which advances the invoice consumes. Conversely, invoicing
before any work can create a receivable and contract liability together.
Both schedules preserve the same earned amounts under these terms. The
finite method retains these billing identities alongside the two performance
calculations. Nonzero impairment, refundable excess payments, contract losses
and tax require their own measurements before this example can be joined to
a more general company book.
